%! AP Statistics — Unit 3, Lesson 3.2 — Homework KEY
\documentclass[10pt]{article}
\usepackage{apstats-article}
\usepackage{apstats-key}
\usepackage{grading}
\gradingassignment{apstats-u03-l02-hw}
\gradingcourse{advanced_statistics}

\begin{document}

\pageheader{Unit 3, Lesson 3.2}{Homework — KEY}
{\small\itshape Write each final answer in its box. Only what is inside a box is graded.}\par

\vspace{0.1in}

\begin{practicebox}
\small
\textbf{1.} Diet and blood pressure.

Study type: \ans{Observational study.} \quad Rand. selection: \ans{Yes.} \quad Rand. assignment: \ans{No.}\\[4pt]
Causation: \ans{No.} \quad Generalize: \ans{Yes (to city adults).}\gradedkey{\mbox{Type~\gradedblank[id=hw1a, type=short, points=1, answer={observational study}, check=text]{1.3in}{0.2in}\enspace Sel.~\gradedblank[id=hw1b, type=short, points=1, answer=yes, check=text]{0.45in}{0.2in}\enspace Assign.~\gradedblank[id=hw1c, type=short, points=1, answer=no, check=text]{0.8in}{0.2in}\enspace Cause~\gradedblank[id=hw1d, type=short, points=1, answer=no, check=text]{0.65in}{0.2in}\enspace Gen.~\gradedblank[id=hw1e, type=short, points=1, answer=yes, check=text]{0.9in}{0.2in}}}\vspace{4pt}
\ans{This is an observational study — the researcher observed but did not impose diets. Confounding variables (exercise, genetics) could explain both diet quality and blood pressure. Random selection allows generalization to the city adult population.}\gradedkeyline{\gradedblank[id=hw1f, type=rubric, points=2, display={Observational: confounding (exercise, genetics), so no causation. Random selection: generalize to city adults.}]{5.4in}{\gradedlinesheight{2}}}

\vspace{6pt}
\textbf{2.} Standing desk experiment.

Study type: \ans{Experiment.} \quad Rand. selection: \ans{No (recruited volunteers).} \quad Rand. assignment: \ans{Yes.}\\[4pt]
Causation: \ans{Yes.} \quad Generalize: \ans{No (only volunteers at this company).}\gradedkey{\mbox{Type~\gradedblank[id=hw2a, type=short, points=1, answer=experiment, check=text]{1.3in}{0.2in}\enspace Sel.~\gradedblank[id=hw2b, type=short, points=1, answer=no, check=text]{0.45in}{0.2in}\enspace Assign.~\gradedblank[id=hw2c, type=short, points=1, answer=yes, check=text]{0.8in}{0.2in}\enspace Cause~\gradedblank[id=hw2d, type=short, points=1, answer=yes, check=text]{1in}{0.2in}\enspace Gen.~\gradedblank[id=hw2e, type=short, points=1, answer=no, check=text]{1in}{0.2in}}}\vspace{4pt}
\ans{Random assignment was used, so the standing desk caused the reduction in back pain. However, employees were not randomly selected, so results cannot generalize beyond similar employees at this company.}\gradedkeyline{\gradedblank[id=hw2f, type=rubric, points=2, display={Random assignment, so causation. Volunteers, not randomly selected: no generalizing beyond this company.}]{5.4in}{\gradedlinesheight{2}}}

\tcbbreak
\vspace{6pt}
\textbf{3.} Reading intervention.
\begin{enumerate}[label=(\alph*), leftmargin=2em, itemsep=5pt]
  \item \ans{Yes — random assignment was used, so the intervention caused the significantly higher scores.}\gradedkeyline{\gradedblank[id=hw3a, type=rubric, points=2, display={Yes: random assignment, so the intervention caused the higher scores.}]{6in}{\gradedlinesheight{2}}}
  \item \ans{No — students were volunteers from a single school, not randomly selected from all third-graders. Results cannot generalize beyond similar volunteers.}\gradedkeyline{\gradedblank[id=hw3b, type=rubric, points=2, display={No: volunteers from one school, not a random sample of all third-graders.}]{6in}{\gradedlinesheight{2}}}
  \item \ans{Randomly select students from a broader population of third-graders (e.g., from a district registry) rather than recruiting volunteers from one school.}\gradedkeyline{\gradedblank[id=hw3c, type=rubric, points=2, display={Randomly select students from all third-graders (e.g., a district registry).}]{6in}{\gradedlinesheight{2}}}
\end{enumerate}
\end{practicebox}

\end{document}
